\documentclass[aspectratio=169, 10pt]{beamer}
\usepackage{beamer_theme_cienciadados2026}

\title[Aula 12 - Matplotlib]{Ciência dos Dados I: Análise Exploratória}
\subtitle{Aula 12: Visualização de Dados com Matplotlib: Interface Orientada a Objetos}
\author{Prof. Dr. Ney Lemke}
\institute{UNESP -- Instituto de Biociências de Botucatu}
\date{Versão 2026}

\begin{document}

\frame{\titlepage}

\begin{frame}{Sumário da Aula}
  \tableofcontents
\end{frame}

\section{Princípios da Comunicação Visual de Dados}

\begin{frame}{A Arte e Ciência dos Gráficos Eficazes}
  Segundo Edward Tufte (\textit{The Visual Display of Quantitative Information}):
  \begin{itemize}
    \item \textbf{Data-Ink Ratio}: Maximize a tinta dedicada aos dados reais; minimize ruído ornamental inútil (\textit{chartjunk}).
    \item \textbf{Integridade Gráfica}: O tamanho visual do efeito exibido na figura deve ser estritamente proporcional ao efeito numérico real dos dados.
    \item \textbf{Contexto e Clareza}: Eixos sempre rotulados com unidades de medida e títulos explicativos.
  \end{itemize}
\end{frame}

\section{A Arquitetura do Matplotlib}

\begin{frame}{Anatomia de uma Figura}
  \begin{columns}
    \begin{column}{0.5\textwidth}
      \begin{block}{Hierarquia de Objetos}
        \begin{itemize}
          \item \textbf{Figure}: O canvas total ou janela que contém um ou mais subplots.
          \item \textbf{Axes}: A área de plotagem propriamente dita (onde ficam os dados, eixos e curvas).
          \item \textbf{Axis}: A escala numérica do eixo X ou Y com seus \textit{ticks} e rótulos.
        \end{itemize}
      \end{block}
    \end{column}
    \begin{column}{0.48\textwidth}
      \begin{alertblock}{Evite \texttt{plt.plot()} solto}
        A interface Orientada a Objetos com \texttt{fig, ax = plt.subplots()} é muito mais explícita, robusta e modular para figuras complexas de múltiplos painéis.
      \end{alertblock}
    \end{column}
  \end{columns}
\end{frame}

\section{Construção Prática e Exportação}

\begin{frame}[fragile]{Estrutura Canônica no Matplotlib}
  \begin{lstlisting}
import matplotlib.pyplot as plt

fig, ax = plt.subplots(figsize=(8, 4.5), dpi=150)

# Plotagem
ax.scatter(df['horsepower'], df['price'], alpha=0.7, color='#1B4973')

# Customizacao explicita
ax.set_title("Potência vs. Preço de Automóveis", fontsize=14, pad=10)
ax.set_xlabel("Potência (HP)")
ax.set_ylabel("Preço (USD)")
ax.grid(True, linestyle='--', alpha=0.5)

# Exportacao de alta qualidade para publicacao
fig.savefig("grafico_final.pdf", bbox_inches='tight')
  \end{lstlisting}
\end{frame}

\begin{frame}{Resumo da Aula}
  \begin{itemize}
    \item Visualização não é cosmética; é ferramenta de descoberta de padrões, tendências e outliers.
    \item A abordagem OO (\texttt{fig, ax}) dá controle total sobre cada elemento gráfico.
    \item \textbf{Prática}: Scatter plots, barras, formatação visual avançada e exportação vetorial.
  \end{itemize}
\end{frame}

\end{document}
